% !TeX program = xelatex
% Bao cao tieu luan - bien soan cho hoc phan Bieu dien tri thuc va ung dung
% Bien dich: XeLaTeX (khuyen nghi dung Overleaf - da co san font Noto Serif va cac goi can thiet)
\documentclass[12pt,a4paper]{report}

% ---------- Ngon ngu va font (XeLaTeX + Unicode tieng Viet) ----------
\usepackage{fontspec}
\usepackage{polyglossia}
\setmainlanguage{vietnamese}
\setmainfont{Times New Roman}
\setsansfont{Arial}
\setmonofont{Consolas}[Scale=0.88]

% ---------- Trang va bo cuc ----------
\usepackage[a4paper,margin=2.7cm]{geometry}
\usepackage{setspace}
\onehalfspacing
\usepackage{titlesec}
\usepackage{fancyhdr}
\setlength{\headheight}{14.5pt}
\pagestyle{fancy}
\fancyhf{}
\fancyhead[L]{\small Biểu diễn tri thức và ứng dụng}
\fancyhead[R]{\small \leftmark}
\fancyfoot[C]{\thepage}
\renewcommand{\headrulewidth}{0.4pt}

% ---------- Bang, hinh, code, sieu lien ket ----------
\usepackage{graphicx}
\usepackage{booktabs}
\usepackage{array}
\usepackage{multirow}
\usepackage{caption}
\usepackage{float}
\usepackage{amsmath,amssymb}
\usepackage{enumitem}
\usepackage{tikz}
\usetikzlibrary{shapes.geometric,arrows.meta,positioning,fit,backgrounds}
\usepackage[most]{tcolorbox}
\usepackage{listings}
\usepackage{xcolor}
\usepackage[colorlinks=true,linkcolor=blue!50!black,citecolor=blue!50!black,urlcolor=blue!50!black]{hyperref}

\definecolor{codebg}{RGB}{245,245,245}
\definecolor{codekw}{RGB}{0,90,160}
\definecolor{codestr}{RGB}{160,20,20}
\definecolor{codecom}{RGB}{110,110,110}

\lstdefinestyle{py}{
  backgroundcolor=\color{codebg},
  basicstyle=\ttfamily\small,
  keywordstyle=\color{codekw}\bfseries,
  stringstyle=\color{codestr},
  commentstyle=\color{codecom}\itshape,
  numbers=left,
  numberstyle=\tiny\color{gray},
  numbersep=6pt,
  frame=single,
  rulecolor=\color{gray!40},
  breaklines=true,
  showstringspaces=false,
  tabsize=4,
  language=Python,
  captionpos=b,
}
\lstdefinestyle{cypher}{
  backgroundcolor=\color{codebg},
  basicstyle=\ttfamily\small,
  numbers=left,
  numberstyle=\tiny\color{gray},
  numbersep=6pt,
  frame=single,
  rulecolor=\color{gray!40},
  breaklines=true,
  showstringspaces=false,
  captionpos=b,
}

\newcommand{\code}[1]{\texttt{\small #1}}

% ==============================================================
\begin{document}

% ------------------------- TRANG BIA -------------------------
\begin{titlepage}
\centering
{\large HỌC VIỆN NGÂN HÀNG\par}
{\large KHOA HỆ THỐNG THÔNG TIN QUẢN LÝ\par}
\vspace{0.3cm}
\rule{0.6\textwidth}{0.4pt}\par
\vspace{2.0cm}

{\Large\bfseries BÁO CÁO TIỂU LUẬN\par}
\vspace{0.4cm}
{\large Học phần: Biểu diễn tri thức và ứng dụng (Cao học)\par}
\vspace{1.6cm}

{\huge\bfseries Hệ thống hỗ trợ phân tích và chẩn đoán \\[0.3em]
sự cố hệ thống mạng và dịch vụ CNTT \\[0.3em]
sử dụng Knowledge Graph, GraphRAG \\[0.3em]
và Rule-Based Reasoning\par}

\vspace{2.2cm}

\begin{flushright}
\begin{tabular}{ll}
Học viên thực hiện: & \textbf{Khang} \\
Giảng viên hướng dẫn: & Giảng viên hướng dẫn học phần \\
\end{tabular}
\end{flushright}

\vfill
{\large \today\par}
\end{titlepage}

\tableofcontents
\clearpage

% ==============================================================
\chapter*{Tóm tắt}
\addcontentsline{toc}{chapter}{Tóm tắt}

Báo cáo trình bày thiết kế, cài đặt và đánh giá thực nghiệm của một hệ thống hỗ trợ chẩn đoán
nguyên nhân gốc rễ (root cause) sự cố trên hệ thống mạng và dịch vụ CNTT, kết hợp ba kỹ thuật
biểu diễn và suy diễn tri thức đã học trong học phần: \textbf{Knowledge Graph (KG)} để biểu diễn
tường minh cấu trúc phụ thuộc giữa thiết bị/dịch vụ, \textbf{suy diễn dựa luật (rule-based
reasoning)} theo mô hình production system cổ điển với cơ chế kết hợp độ tin cậy kiểu MYCIN
(certainty factor), và \textbf{GraphRAG} (Retrieval-Augmented Generation trên đồ thị tri thức)
để khai thác một mô hình ngôn ngữ lớn (LLM) sinh giải thích dựa trên ngữ cảnh được xác thực từ
đồ thị thay vì suy đoán tự do.

Hệ thống được cài đặt đầy đủ bằng Python, dùng Neo4j làm hạ tầng lưu trữ đồ thị tri thức, và được
đánh giá trên hai tập dữ liệu độc lập: (1) một tập dữ liệu mô phỏng mạng viễn thông gồm 9 thiết bị,
2 site và 5 kịch bản sự cố mẫu, đúng phạm vi đã đăng ký trong đề cương; và (2) một tập
\textbf{dữ liệu thật} gồm 78 sự cố hạ tầng IT-ops/microservice thật, lấy từ bộ dữ liệu công khai
NetManAIOps/DejaVu, dùng để kiểm chứng độc lập khả năng tổng quát hoá của phương pháp trên một hệ
thống có cấu trúc phụ thuộc thực tế mà nhóm không tự tạo ra. Trên cả hai tập dữ liệu, cả ba
phương án --- chỉ dùng luật, chỉ dùng GraphRAG, và kết hợp cả hai --- đều đạt độ chính xác top-1
100\% với LLM mô phỏng (MockLLM). Khi thay bằng một LLM thật (Groq/Llama~3.3-70B), báo cáo ghi
nhận và phân tích một phát hiện thực nghiệm quan trọng: LLM thật có thể tự tin nhưng sai, và module
tổng hợp cần được thiết kế lại để không cho một câu trả lời GraphRAG không được luật nào xác nhận
lấn át một câu trả lời đúng từ Rule Engine --- đúng minh hoạ cho vai trò "kiểm chứng chéo" của suy
diễn dựa luật mà lý thuyết của học phần đề cập.

% ==============================================================
\chapter{Giới thiệu}

\section{Lý do chọn đề tài}

Trong vận hành mạng viễn thông, kỹ sư trung tâm điều hành mạng (NOC -- Network Operations Center)
hằng ngày phải xử lý một khối lượng lớn cảnh báo (alarm), chỉ số hiệu năng (KPI) và log thiết bị
đến từ nhiều lớp mạng khác nhau (truyền dẫn, truy nhập, core, dịch vụ). Việc xác định nguyên nhân
gốc rễ của sự cố hiện chủ yếu dựa vào kinh nghiệm cá nhân và tra cứu tài liệu rời rạc, dẫn đến thời
gian xử lý sự cố (MTTR) kéo dài và khó chia sẻ tri thức giữa các kỹ sư.

Bài toán này phù hợp để áp dụng các kỹ thuật biểu diễn tri thức đã học: tri thức về cấu trúc mạng,
quan hệ phụ thuộc giữa thiết bị và dịch vụ, cũng như kinh nghiệm xử lý sự cố có thể được mô hình
hoá tường minh dưới dạng Knowledge Graph; việc suy luận nguyên nhân có thể được hỗ trợ bởi cả suy
diễn dựa luật và truy xuất ngữ cảnh bằng GraphRAG kết hợp LLM. Đây đồng thời là một bài toán kinh
điển trong lĩnh vực AIOps (AI for IT Operations), nên ngoài mạng viễn thông, phương pháp còn được
kiểm chứng thêm trên dữ liệu hạ tầng CNTT/microservice thật (Mục~\ref{ch:data}), cho thấy tính tổng
quát của cách tiếp cận.

\section{Mục tiêu đề tài}

\textbf{Mục tiêu tổng quát:} Xây dựng hệ thống hỗ trợ kỹ sư vận hành phân tích, chẩn đoán nguyên
nhân gốc rễ của sự cố và đề xuất hướng xử lý dựa trên tri thức được biểu diễn dưới dạng đồ thị tri
thức.

\textbf{Mục tiêu cụ thể:}
\begin{enumerate}[label=(\arabic*)]
    \item Thiết kế mô hình Knowledge Graph biểu diễn thiết bị, cảnh báo, KPI và các mối quan hệ
    (kết nối vật lý/logic, phụ thuộc dịch vụ, nguyên nhân -- hệ quả).
    \item Xây dựng bộ luật suy diễn (rule-based reasoning) mã hoá kinh nghiệm chẩn đoán sự cố của
    kỹ sư vận hành dưới dạng luật IF-THEN (production rules).
    \item Xây dựng pipeline GraphRAG: khi có sự cố, truy xuất các thực thể và quan hệ liên quan
    trong KG, kết hợp ngữ cảnh này với LLM để sinh giải thích và đề xuất nguyên nhân.
    \item Kết hợp kết quả suy diễn luật và kết quả truy xuất GraphRAG thành khuyến nghị nguyên
    nhân gốc rễ và hướng xử lý cho kỹ sư vận hành, với cơ chế kiểm chứng chéo.
    \item Xây dựng demo/prototype minh hoạ toàn bộ luồng xử lý, gồm cả giao diện web, trên một tập
    kịch bản sự cố mẫu và một tập dữ liệu sự cố thật độc lập.
\end{enumerate}

\section{Đối tượng và phạm vi nghiên cứu}

\textbf{Đối tượng nghiên cứu:} các phương pháp biểu diễn tri thức bằng đồ thị (Knowledge Graph),
kỹ thuật Retrieval-Augmented Generation trên đồ thị (GraphRAG), và các phương pháp suy diễn dựa
luật (rule-based/production system).

\textbf{Phạm vi dữ liệu} gồm hai lớp:
\begin{itemize}
    \item \emph{Lớp chính:} dữ liệu mô phỏng/giả lập về topology mạng viễn thông, cảnh báo và KPI
    ở quy mô vừa phải (một số loại thiết bị tiêu biểu: BTS/gNodeB, OLT, router, thiết bị truyền
    dẫn), do giới hạn thời gian của học phần không đi sâu vào việc kết nối dữ liệu vận hành thực
    tế đang bảo mật.
    \item \emph{Lớp đối chiếu bổ sung:} dữ liệu sự cố \textbf{thật} ngoài miền viễn thông (hạ tầng
    IT-ops/microservice, bộ dữ liệu công khai NetManAIOps/DejaVu, giấy phép CC-BY~4.0), dùng để
    kiểm chứng độc lập khả năng tổng quát hoá của phương pháp trên một hệ thống có cấu trúc phụ
    thuộc thực tế, không tự tạo dữ liệu (xem Mục~\ref{ch:data}).
\end{itemize}

\textbf{Phạm vi kỹ thuật:} tập trung vào việc mô hình hoá tri thức, xây dựng bộ luật suy diễn và
luồng GraphRAG kết hợp LLM; không đi sâu vào tối ưu hạ tầng triển khai production.

% ==============================================================
\chapter{Cơ sở lý thuyết}

\section{Biểu diễn tri thức bằng Knowledge Graph}

Tri thức về hệ thống mạng được mô hình hoá dưới dạng đồ thị có hướng gồm các thực thể (node) và
quan hệ (edge). Mô hình dữ liệu được tổ chức theo hướng ontology đơn giản (class -- property --
instance), thuận lợi cho việc truy vấn quan hệ nhiều bước (multi-hop) khi chẩn đoán sự cố lan
truyền qua nhiều lớp mạng. Sáu lớp thực thể được định nghĩa: \emph{Device}, \emph{Alarm},
\emph{KPI}, \emph{Service}, \emph{Site}, \emph{Incident}; và tám loại quan hệ:
\code{CONNECTS\_TO}, \code{DEPENDS\_ON}, \code{BELONGS\_TO\_SITE}, \code{RAISED\_ALARM},
\code{AFFECTS\_SERVICE}, \code{CAUSED\_BY}, \code{HAS\_KPI}, \code{PART\_OF\_INCIDENT}.

Một điểm quan trọng được rút ra trong quá trình cài đặt (Mục~\ref{sec:daghoa}): quan hệ
\code{DEPENDS\_ON} không nên bị ràng buộc thành một cây đơn-cha (mỗi thiết bị chỉ có một cha), mà
cần biểu diễn như một \textbf{đồ thị có hướng không chu trình (DAG)} cho phép nhiều-cha, vì đây mới
là cấu trúc phụ thuộc thực tế của cả mạng viễn thông (dự phòng, dual-homing) lẫn hạ tầng
microservice (một dịch vụ gọi nhiều backend và đồng thời phụ thuộc vào máy chủ vật lý nó chạy
trên).

\section{Suy diễn dựa luật (Rule-Based Reasoning)}

Suy diễn dựa luật được cài đặt theo mô hình \emph{production system} kinh điển: mỗi luật có dạng
\textbf{IF (điều kiện trên các sự kiện quan sát được trong KG) THEN (giả thuyết nguyên nhân, kèm
độ tin cậy)}. Khi nhiều luật cùng chỉ ra một thiết bị, độ tin cậy được kết hợp theo công thức
\emph{certainty factor} của hệ chuyên gia MYCIN~\cite{shortliffe1975mycin}:
\begin{equation}
    CF_{1,2} = CF_1 + CF_2 \times (1 - CF_1)
    \label{eq:cf}
\end{equation}
Công thức~(\ref{eq:cf}) có hai tính chất phù hợp với bài toán: (i) đơn điệu tăng --- càng nhiều
bằng chứng độc lập ủng hộ cùng một giả thuyết thì độ tin cậy tổng hợp càng cao; (ii) bị chặn trong
$[0,1]$, tránh hiện tượng độ tin cậy vượt ngưỡng khi kết hợp nhiều nguồn.

\section{GraphRAG}

GraphRAG (Retrieval-Augmented Generation trên đồ thị tri thức) là kỹ thuật kết hợp truy xuất thông
tin có cấu trúc từ một Knowledge Graph với khả năng sinh văn bản của LLM, nhằm hạn chế hiện tượng
LLM "ảo giác" (hallucination) khi suy luận trên tri thức chuyên biệt. Quy trình gồm ba bước: (1)
xác định các thực thể liên quan đến sự cố (thiết bị, cảnh báo, KPI bất thường) --- gọi là
\emph{seed entities}; (2) truy xuất tiểu đồ thị (subgraph) đa bước quanh các seed entities này,
bao gồm cả các sự cố lịch sử tương tự; (3) chuyển tiểu đồ thị và ngữ cảnh liên quan thành văn bản
có cấu trúc để đưa vào LLM, giúp LLM sinh giải thích/nguyên nhân dựa trên tri thức đã được xác thực
thay vì suy đoán tự do.

\section{Kết hợp luật và GraphRAG: kiểm chứng chéo}
\label{sec:crossval-theory}

Vì rule-based reasoning dựa trên tri thức tường minh, đã kiểm chứng, còn GraphRAG dựa trên suy
luận xác suất của một mô hình thống kê (LLM), hai phương pháp bổ trợ lẫn nhau: luật thu hẹp không
gian nguyên nhân một cách chắc chắn, còn LLM có khả năng diễn giải các mẫu hình chưa được luật hoá
tường minh (ví dụ so khớp với sự cố lịch sử tương tự). Mục~\ref{ch:ket-qua} trình bày một phát hiện
thực nghiệm cho thấy vai trò "kiểm chứng chéo" này không chỉ là lý thuyết suông mà thực sự cần
thiết khi làm việc với LLM thật.

% ==============================================================
\chapter{Thiết kế và kiến trúc hệ thống}

\section{Kiến trúc tổng thể}

Hệ thống gồm bốn thành phần chính, tương ứng bốn package Python (\code{src/kg}, \code{src/rules},
\code{src/graphrag}, \code{src/synthesis}), giao tiếp qua một cơ sở dữ liệu đồ thị Neo4j dùng
chung (Hình~\ref{fig:kientruc}):

\begin{enumerate}[label=(\arabic*)]
    \item \textbf{Lớp dữ liệu và Knowledge Graph} (\code{src/kg}): định nghĩa schema, nạp dữ liệu
    thiết bị/cảnh báo/KPI/quan hệ mạng vào Neo4j.
    \item \textbf{Module suy diễn luật -- Rule Engine} (\code{src/rules}): áp dụng các luật
    IF-THEN để đề xuất tập nguyên nhân ứng viên ban đầu.
    \item \textbf{Module GraphRAG} (\code{src/graphrag}): truy xuất tiểu đồ thị liên quan, sinh
    giải thích và đề xuất nguyên nhân bằng LLM dựa trên ngữ cảnh truy xuất được.
    \item \textbf{Module tổng hợp -- Synthesis} (\code{src/synthesis}): kết hợp kết quả từ hai
    module trên, xếp hạng nguyên nhân theo độ tin cậy, đề xuất hướng xử lý.
\end{enumerate}

Bên trên bốn thành phần lõi là một giao diện web (\code{app.py}, Chương~\ref{ch:demo}) và hai kịch
bản dùng lại toàn bộ pipeline nhưng trên hai nguồn dữ liệu khác nhau: kịch bản demo/đánh giá viễn
thông (\code{src/demo.py}, \code{src/evaluation/evaluate.py}) và kịch bản đánh giá trên dữ liệu
thật DejaVu (\code{src/evaluation/evaluate\_dejavu.py}).

\begin{figure}[H]
\centering
\begin{tikzpicture}[
    box/.style={rectangle, rounded corners, draw=black!60, fill=blue!8, minimum width=3.6cm,
                minimum height=1.15cm, align=center, font=\small},
    dbbox/.style={cylinder, shape border rotate=90, draw=black!60, fill=orange!15,
                  minimum width=2.6cm, minimum height=1.7cm, align=center, font=\small},
    arrow/.style={-{Latex[length=2mm]}, thick}
]
\node[box]                            (kg)      at (-4.2, 5.0) {Lớp dữ liệu \& KG\\\texttt{src/kg}};
\node[box]                            (rules)   at ( 4.2, 5.0) {Rule Engine\\\texttt{src/rules}};
\node[dbbox]                          (neo4j)   at ( 0.0, 3.0) {Neo4j\\Knowledge Graph};
\node[box]                            (graphrag) at (-4.2, 1.0) {Module GraphRAG\\\texttt{src/graphrag}};
\node[box]                            (synth)   at ( 4.2, 1.0) {Module tổng hợp\\\texttt{src/synthesis}};
\node[box, fill=green!10, minimum width=5.6cm] (ui) at (0.0, -1.6) {Giao diện web (Streamlit)\\\texttt{app.py}};

\draw[arrow] (kg) -- (neo4j);
\draw[arrow] (neo4j) -- (rules);
\draw[arrow] (neo4j) -- (graphrag);
\draw[arrow] (rules) -- (synth);
\draw[arrow] (graphrag) -- (synth);
\draw[arrow] (graphrag.south) |- (ui.west);
\draw[arrow] (synth.south) |- (ui.east);
\end{tikzpicture}
\caption{Kiến trúc tổng thể của hệ thống: bốn module lõi giao tiếp qua Neo4j, tổng hợp kết quả
lên giao diện web.}
\label{fig:kientruc}
\end{figure}

\section{Schema Knowledge Graph}

Bảng~\ref{tab:schema} tóm tắt schema KG. Toàn bộ ràng buộc \code{UNIQUE} trên thuộc tính
\code{id} của mỗi nhãn được khai báo bằng \emph{constraint} của Neo4j, đảm bảo tính toàn vẹn khi
nạp lại dữ liệu nhiều lần.

\begin{table}[H]
\centering
\caption{Schema Knowledge Graph}
\label{tab:schema}
\begin{tabular}{@{}p{3.2cm}p{9.8cm}@{}}
\toprule
\textbf{Thành phần} & \textbf{Mô tả} \\
\midrule
\multicolumn{2}{@{}l}{\textit{Nhãn thực thể (node label)}} \\
\code{Device} & Thiết bị mạng/hạ tầng (transmission, router, olt, gnodeb, docker, db, os, osb...) \\
\code{Alarm} & Một cảnh báo cụ thể, gắn với thiết bị phát sinh (\code{type}, \code{severity}) \\
\code{KPI} & Một chỉ số đo lường bất thường (\code{name}, \code{value}, \code{status}) \\
\code{Service} & Dịch vụ nghiệp vụ bị ảnh hưởng (Mobile Data, Broadband...) \\
\code{Site} & Vị trí vật lý của thiết bị \\
\code{Incident} & Một sự cố (đang xảy ra hoặc lịch sử), gộp các Alarm/KPI liên quan \\
\midrule
\multicolumn{2}{@{}l}{\textit{Quan hệ (relationship type)}} \\
\code{DEPENDS\_ON} & Thiết bị con phụ thuộc thiết bị cha (đồ thị DAG, cho phép nhiều cha) \\
\code{CONNECTS\_TO} & Kết nối vật lý/logic giữa hai thiết bị \\
\code{RAISED\_ALARM} & Thiết bị $\rightarrow$ cảnh báo mà nó phát sinh \\
\code{HAS\_KPI} & Thiết bị $\rightarrow$ chỉ số KPI đo được \\
\code{AFFECTS\_SERVICE} & Thiết bị $\rightarrow$ dịch vụ bị ảnh hưởng nếu thiết bị lỗi \\
\code{BELONGS\_TO\_SITE} & Thiết bị $\rightarrow$ site vật lý \\
\code{PART\_OF\_INCIDENT} & Alarm/KPI $\rightarrow$ sự cố đang xét \\
\code{CAUSED\_BY} & Incident $\rightarrow$ thiết bị là nguyên nhân gốc rễ (ground truth) \\
\bottomrule
\end{tabular}
\end{table}

\section{Tổng quát hoá sang đồ thị phụ thuộc đa-cha (DAG)}
\label{sec:daghoa}

Thiết kế ban đầu giả định mỗi thiết bị chỉ phụ thuộc duy nhất một thiết bị cha (\code{depends\_on}
là một giá trị đơn), phù hợp với topology hình cây của tập dữ liệu viễn thông mô phỏng. Khi tích
hợp dữ liệu thật DejaVu (Chương~\ref{ch:data}), topology thực tế cho thấy một container có thể
đồng thời gọi nhiều container/database khác \emph{và} phụ thuộc vào máy chủ vật lý nó chạy trên
--- tức là \textbf{nhiều-cha}. Toàn bộ Rule Engine, module truy xuất subgraph của GraphRAG và
module suy diễn của MockLLM đã được tổng quát hoá: \code{parent\_id} (một giá trị) được thay bằng
\code{parents\_of[device]} (một danh sách), các luật liên quan đến "thiết bị cha" được viết lại để
kiểm tra \emph{có tồn tại một cha nào bị cảnh báo hay không} thay vì so sánh với một cha duy nhất.
Đây là một tổng quát hoá hợp lý ngay cả với mạng viễn thông thật (một router có thể có hai đường
uplink dự phòng), không phải một giải pháp vá lỗi riêng cho dữ liệu DejaVu. Toàn bộ 5 kịch bản viễn
thông vẫn cho kết quả đúng 100\% sau khi tổng quát hoá (Bảng~\ref{tab:eval-telecom}), xác nhận thay
đổi không phá vỡ hành vi cũ.

% ==============================================================
\chapter{Cài đặt chi tiết}

\section{Bộ luật suy diễn (Rule Engine)}

Bốn luật được xây dựng, mã hoá các mẫu hình chẩn đoán phổ biến trong vận hành NOC
(Bảng~\ref{tab:rules}). Việc lựa chọn bốn mẫu hình này dựa trên phân tích cách một kỹ sư NOC thường
suy luận: từ mẫu hình lan truyền theo cấu trúc (luật R1, R3) đến mẫu hình khu trú
(luật R2) và mẫu hình cảnh báo sớm dựa trên xu hướng KPI (luật R4).

\begin{table}[H]
\centering
\caption{Bốn luật IF-THEN của Rule Engine}
\label{tab:rules}
\begin{tabular}{@{}p{2.6cm}p{7.6cm}p{1.6cm}@{}}
\toprule
\textbf{Luật} & \textbf{Điều kiện (IF) $\rightarrow$ Kết luận (THEN)} & \textbf{CF} \\
\midrule
R1: Cascading parent failure &
IF thiết bị $X$ phát cảnh báo mất kết nối (Link\_Down) VÀ toàn bộ thiết bị con của $X$ đều phát
cảnh báo mất tín hiệu THEN nguyên nhân gốc rễ là $X$ & 0.95 \\[4pt]
R2: Isolated device fault &
IF thiết bị $X$ phát cảnh báo VÀ không có thiết bị cha nào bị cảnh báo VÀ không có thiết bị con nào
bị cảnh báo THEN nguyên nhân gốc rễ là $X$ & 0.85 \\[4pt]
R3: Shared-parent sibling pattern &
IF $\geq 2$ thiết bị con cùng một cha $P$ đều phát cảnh báo VÀ $P$ chưa có cảnh báo tường minh THEN
$P$ là điểm lỗi chung & 0.75 \\[4pt]
R4: KPI-only early warning &
IF thiết bị $X$ có KPI vượt ngưỡng cảnh báo/nguy hiểm VÀ chưa có alarm chính thức THEN nghi ngờ $X$
đang suy giảm hiệu năng & 0.50--0.65 \\
\bottomrule
\end{tabular}
\end{table}

Luật R1 chính là ví dụ đã nêu trong đề cương ban đầu ("NẾU thiết bị truyền dẫn X mất kết nối VÀ
tất cả thiết bị con của X đều phát cảnh báo mất tín hiệu THÌ nguyên nhân gốc rễ nhiều khả năng là
sự cố tại X"). Khi nhiều luật cùng chỉ ra một thiết bị, độ tin cậy được kết hợp bằng công
thức~(\ref{eq:cf}) (Đoạn mã~\ref{lst:cf}).

\begin{lstlisting}[style=py,caption={Kết hợp độ tin cậy kiểu MYCIN (\texttt{src/rules/rule\_engine.py})},label=lst:cf]
def _combine_certainty_factors(confidences):
    """MYCIN-style certainty-factor combination: CF12 = CF1 + CF2*(1-CF1)."""
    combined = 0.0
    for c in confidences:
        combined = combined + c * (1 - combined)
    return round(combined, 4)
\end{lstlisting}

Đoạn mã~\ref{lst:ruleA} minh hoạ cách cài đặt luật R1 bằng Python thuần, thao tác trên các cấu trúc
dữ liệu (\code{devices}, \code{children\_of}, \code{alarms\_by\_device}) được truy vấn một lần từ
Neo4j (hàm \code{fetch\_facts}) rồi đánh giá lại nhiều lần trong bộ nhớ, tránh việc phải viết logic
suy diễn phức tạp trực tiếp bằng Cypher.

\begin{lstlisting}[style=py,caption={Luật R1 -- cascading parent failure (\texttt{src/rules/rule\_engine.py})},label=lst:ruleA]
def rule_cascading_parent_failure(facts):
    candidates = []
    devices, children_of = facts["devices"], facts["children_of"]
    alarms_by_device = facts["alarms_by_device"]
    for dev_id in devices:
        if not _has_alarm_of_type(alarms_by_device, dev_id, LINK_FAILURE_ALARM_TYPES):
            continue
        children = children_of.get(dev_id, [])
        if not children:
            continue
        if all(_has_alarm_of_type(alarms_by_device, c, SIGNAL_LOSS_ALARM_TYPES)
               for c in children):
            candidates.append(RuleCandidate(
                device_id=dev_id, confidence=0.95,
                rule_id="R1_cascading_parent_failure", ...))
    return candidates
\end{lstlisting}

\section{Pipeline GraphRAG}

Pipeline GraphRAG (\code{src/graphrag/pipeline.py}) gồm ba bước tuần tự, cài đặt tương ứng trong ba
mô-đun:

\begin{enumerate}[label=(\arabic*)]
    \item \textbf{Retriever} (\code{retriever.py}): xác định \emph{seed entities} (thiết bị có
    alarm/KPI thuộc sự cố đang xét), sau đó truy vấn Cypher để lấy tiểu đồ thị trong bán kính 2
    bước nhảy (\code{DEPENDS\_ON*0..2}) quanh các seed entities, gồm cả alarm/KPI/service/sự cố
    lịch sử liên quan.
    \item \textbf{Context builder} (\code{context\_builder.py}): chuyển tiểu đồ thị (dạng cấu trúc
    Python) thành văn bản tiếng Việt có cấu trúc rõ ràng (theo từng mục: thiết bị liên quan, cảnh
    báo hiện tại, KPI bất thường, dịch vụ bị ảnh hưởng, sự cố lịch sử tương tự).
    \item \textbf{LLM client} (\code{llm\_client.py}): gửi văn bản ngữ cảnh cho một LLM và yêu cầu
    trả về JSON gồm \code{root\_cause\_device}, \code{confidence}, \code{explanation}.
\end{enumerate}

Đoạn mã~\ref{lst:cypher} là truy vấn Cypher truy xuất tiểu đồ thị đa bước --- điểm mấu chốt để
GraphRAG có thể suy luận trên các mẫu hình lan truyền nhiều tầng thay vì chỉ nhìn một thiết bị đơn
lẻ.

\begin{lstlisting}[style=cypher,caption={Truy vấn Cypher truy xuất tiểu đồ thị đa bước},label=lst:cypher]
MATCH (seed:Device) WHERE seed.id IN $seed_ids
MATCH (seed)-[:DEPENDS_ON*0..2]-(neighbor:Device)
OPTIONAL MATCH (neighbor)-[:DEPENDS_ON]->(parent:Device)
OPTIONAL MATCH (neighbor)-[:BELONGS_TO_SITE]->(site:Site)
RETURN neighbor.id AS id, neighbor.type AS type, neighbor.name AS name,
       parent.id AS parent_id, site.name AS site_name
\end{lstlisting}

\textbf{Thiết kế pluggable cho LLM.} Lớp \code{LLMClient} trừu tượng được cài đặt bởi bốn lớp con:
\code{MockLLM} (mặc định, không cần API key, suy luận offline bằng cách so khớp theo hệ số
Jaccard với sự cố lịch sử kết hợp suy luận cấu trúc trên đồ thị), \code{AnthropicLLM},
\code{OpenAILLM} và \code{GroqLLM}. Việc chuyển đổi giữa các nhà cung cấp chỉ cần đổi biến môi
trường \code{LLM\_PROVIDER}, không cần sửa mã nguồn pipeline --- một ví dụ áp dụng nguyên lý tách
biệt giao diện (interface) khỏi cài đặt cụ thể.

Trong quá trình tích hợp LLM thật (Groq, mô hình \code{llama-3.3-70b-versatile}), nhóm phát hiện
LLM thường bọc câu trả lời JSON trong đoạn văn bản suy luận và/hoặc khối markdown \verb|```json|
thay vì trả JSON thuần --- khác với giả định ban đầu. Bộ phân tích JSON được viết lại theo hướng
bền vững hơn: tìm khối \verb|```json ... ```| bất kỳ vị trí nào trong văn bản, nếu không có thì
tìm cặp dấu ngoặc \code{\{...\}} ngoài cùng, mới thử phân tích toàn bộ chuỗi.

\section{Module tổng hợp và cơ chế kiểm chứng chéo}
\label{sec:combiner}

Module \code{src/synthesis/combiner.py} nhận đầu vào là danh sách ứng viên từ Rule Engine và một
kết quả duy nhất từ GraphRAG, gộp theo thiết bị, rồi:

\begin{itemize}
    \item Nếu một thiết bị được \emph{cả hai} nguồn cùng đề xuất (đồng thuận), độ tin cậy được kết
    hợp bằng công thức~(\ref{eq:cf}) rồi cộng thêm một hệ số thưởng
    \code{CROSS\_VALIDATION\_BOOST = 0.1} (giới hạn tối đa 0.99) --- thể hiện tinh thần "hai nguồn
    độc lập cùng đồng ý thì đáng tin hơn một nguồn".
    \item Nếu một thiết bị chỉ được GraphRAG đề xuất mà không có luật nào xác nhận, độ tin cậy tự
    báo cáo của LLM bị \textbf{chiết khấu} bằng hệ số
    \code{UNCORROBORATED\_GRAPHRAG\_DISCOUNT = 0.7} trước khi xếp hạng.
\end{itemize}

Hệ số chiết khấu ở gạch đầu dòng thứ hai \textbf{không} có trong thiết kế ban đầu, mà được thêm vào
sau một phát hiện thực nghiệm thật khi thử nghiệm với Groq (trình bày chi tiết ở
Mục~\ref{sec:llm-thuc-nghiem}): một LLM thật có thể tự báo cáo độ tin cậy cao dù trả lời sai, và
nếu không chiết khấu, câu trả lời sai đó có thể lấn át một ứng viên đúng nhưng chỉ có Rule Engine
ủng hộ. Đây là một minh chứng thực nghiệm cụ thể cho luận điểm lý thuyết ở
Mục~\ref{sec:crossval-theory}: suy diễn dựa luật, vì là tri thức đã được kiểm chứng một cách tường
minh, cần được ưu tiên hơn một phỏng đoán đơn lẻ, không được xác nhận, của một mô hình thống kê.

\begin{lstlisting}[style=py,caption={Chiết khấu ứng viên GraphRAG không được luật nào xác nhận (\texttt{src/synthesis/combiner.py})},label=lst:discount]
effective_graphrag_conf = graphrag_conf
if graphrag_conf is not None and rule_conf is None:
    effective_graphrag_conf = graphrag_conf * UNCORROBORATED_GRAPHRAG_DISCOUNT

confs = [c for c in (rule_conf, effective_graphrag_conf) if c is not None]
combined_conf = _combine_certainty_factors(confs)
if agreement:
    combined_conf = min(MAX_CONFIDENCE, combined_conf + CROSS_VALIDATION_BOOST)
\end{lstlisting}

Cuối cùng, hàm \code{suggest\_action} tra một bảng mẫu câu hướng xử lý theo loại thiết bị
(\code{transmission}, \code{router}, \code{olt}, \code{gnodeb} cho viễn thông; \code{docker},
\code{db}, \code{os}, \code{osb} cho hạ tầng DejaVu), sinh khuyến nghị cụ thể thay vì một thông
điệp chung chung.

% ==============================================================
\chapter{Dữ liệu: hành trình tìm và tích hợp dữ liệu thật}
\label{ch:data}

\section{Vì sao cần tìm dữ liệu thật}

Đề cương ban đầu giới hạn phạm vi ở dữ liệu mô phỏng, vì dữ liệu vận hành thật của các nhà mạng
luôn ở diện bảo mật. Tuy nhiên, để tăng tính thuyết phục của phương pháp, nhóm đã thử tìm các
nguồn dữ liệu công khai, thật, có thể dùng để kiểm chứng độc lập -- không phải để thay thế tập dữ
liệu viễn thông mô phỏng, mà để trả lời câu hỏi: \emph{phương pháp có generalize được ra ngoài dữ
liệu tự tạo hay không?}

\section{Ba nguồn được khảo sát}

Bảng~\ref{tab:datasources} tóm tắt ba nguồn dữ liệu thật đã khảo sát và kết quả xác minh thực tế
(không chỉ dựa trên tên/mô tả, mà đã thử tải trực tiếp).

\begin{table}[H]
\centering
\caption{Ba nguồn dữ liệu thật đã khảo sát}
\label{tab:datasources}
\begin{tabular}{@{}p{2.4cm}p{5.6cm}p{4.2cm}@{}}
\toprule
\textbf{Nguồn} & \textbf{Mô tả} & \textbf{Kết quả xác minh} \\
\midrule
NetRCA~\cite{netrca2022} & Alarm/KPI mạng 5G thật kèm causal graph, từ ICASSP 2022 AIOps
Challenge -- khớp nhất với đúng bài toán & \textbf{Không dùng được}: trang challenge không phân
giải được DNS, không tìm thấy bản sao dữ liệu ở nơi khác \\[4pt]
TeleLogs~\cite{telelogs2025} & Dữ liệu RF tổng hợp mô phỏng UE lái xe qua vùng phủ 5G, host trên
Hugging Face (\texttt{netop/TeleLogs}) & \textbf{Không dùng được}: xác nhận qua API là
\texttt{gated: manual} (phải xin quyền, chờ duyệt thủ công) và chính dataset card cấm
redistribute \\[4pt]
NetManAIOps/DejaVu~\cite{dejavu2022} & Dữ liệu lỗi hạ tầng IT-ops/microservice thật (Docker/DB/OS),
kèm đồ thị phụ thuộc thật và ground-truth root cause, từ Zenodo & \textbf{Dùng được}: tải trực
tiếp không cần đăng nhập, giấy phép CC-BY~4.0 cho phép redistribute một tập con \\
\bottomrule
\end{tabular}
\end{table}

Hành trình này bản thân nó là một minh hoạ thực tế cho một vấn đề thường gặp khi làm việc với dữ
liệu AIOps/viễn thông: dữ liệu vận hành thật gần như luôn bị giới hạn bởi hạ tầng không ổn định
hoặc rào cản pháp lý/quyền truy cập, khiến các nghiên cứu buộc phải dựa vào dữ liệu mô phỏng hoặc
dữ liệu thật từ miền lân cận.

\section{Dữ liệu NetManAIOps/DejaVu}

Bộ dữ liệu được công bố kèm bài báo tại ESEC/FSE~2022~\cite{dejavu2022} (Zenodo, DOI
10.5281/zenodo.6955909, giấy phép CC-BY~4.0). Nhóm viết một script tái tạo được
(\code{scripts/fetch\_dejavu.py}) để tải và phân tích hai tệp gốc:

\begin{itemize}
    \item \code{graph.yml}: mô tả đồ thị phụ thuộc thật gồm 17 thiết bị -- 3 database (\code{db}),
    8 container (\code{docker}), 4 máy chủ vật lý (\code{os}) và 2 nút Online Service Bus
    (\code{osb}) -- với 24 quan hệ gọi dịch vụ (\emph{call edges}) và 8 quan hệ triển khai
    (container chạy trên máy chủ nào). Đây chính là cấu trúc DAG đa-cha đã thúc đẩy tổng quát hoá
    ở Mục~\ref{sec:daghoa}.
    \item \code{faults.csv}: 78 bản ghi sự cố thật, mỗi bản ghi gồm loại lỗi (CPU fault, network
    delay, network loss, db connection limit, db close), thiết bị bị ảnh hưởng và
    \textbf{ground-truth root cause}. Toàn bộ 78 sự cố đều là lỗi đơn-thiết-bị (nguyên nhân luôn
    chính là thiết bị phát sinh lỗi), không có trường hợp lan truyền qua nhiều thiết bị -- khác
    với 5 kịch bản viễn thông vốn được thiết kế riêng để kiểm tra suy luận lan truyền (cascade).
\end{itemize}

Vì vậy, 78 sự cố thật của DejaVu chủ yếu kiểm tra khả năng nhận diện đúng lỗi khu trú
(luật~R2) trên một topology thật, phức tạp, không phải khả năng suy luận cascade -- điều này được
nêu rõ và không bị "làm mờ" trong báo cáo đánh giá (Chương~\ref{ch:ket-qua}).

\textbf{Lưu ý về tính trung thực khoa học:} DejaVu là dữ liệu hạ tầng IT-ops/microservice thật,
\textbf{không phải} dữ liệu viễn thông. Nó được dùng vì có cùng bản chất đồ thị-lý-thuyết (đồ thị
phụ thuộc + lan truyền lỗi + định vị nguyên nhân gốc) với bài toán viễn thông, nên là một phép kiểm
tra tổng quát hoá hợp lệ của \emph{phương pháp}, chứ không phải một minh chứng về hiệu năng trên dữ
liệu viễn thông thật.

% ==============================================================
\chapter{Giao diện demo}
\label{ch:demo}

Để đáp ứng yêu cầu "giao diện demo dạng web app đơn giản" của đề cương, nhóm xây dựng một ứng dụng
web bằng Streamlit (\code{app.py}), gồm hai tab:

\begin{itemize}
    \item \textbf{Tab "Chẩn đoán trực tiếp":} người dùng chọn một trong hai bộ dữ liệu (viễn thông
    mô phỏng hoặc DejaVu thật), chọn một sự cố cụ thể, bấm nút để chạy toàn bộ pipeline. Kết quả
    hiển thị gồm: bảng ứng viên từ Rule Engine kèm luật kích hoạt, chẩn đoán của GraphRAG kèm độ
    tin cậy và giải thích, \textbf{sơ đồ đồ thị con trực quan} (thiết bị là node, quan hệ
    \code{DEPENDS\_ON} là cạnh có hướng, tô màu theo trạng thái: đỏ = đang có alarm, cam = chỉ có
    KPI bất thường, vàng/ngôi sao = nguyên nhân gốc rễ được chẩn đoán), và khuyến nghị tổng hợp
    cuối cùng kèm hướng xử lý cụ thể.
    \item \textbf{Tab "Đánh giá tổng thể":} chạy đánh giá hàng loạt trên một trong hai bộ dữ liệu,
    hiển thị biểu đồ cột so sánh độ chính xác giữa ba phương án (rule-only / GraphRAG-only /
    combined) cùng bảng chi tiết từng trường hợp.
\end{itemize}

Sơ đồ đồ thị con được cài đặt bằng thư viện \code{pyvis} (dựng trên vis-network.js), toàn bộ mã
JavaScript được nhúng trực tiếp vào HTML (\code{cdn\_resources="in\_line"}) để nhúng an toàn vào
iframe của Streamlit mà không phụ thuộc tệp tĩnh bên ngoài.

% ==============================================================
\chapter{Kết quả thực nghiệm và đánh giá}
\label{ch:ket-qua}

\section{Phương pháp đánh giá}

Theo đúng yêu cầu của đề cương, hệ thống được đánh giá bằng cách so sánh độ chính xác top-1 (thiết
bị được xếp hạng cao nhất có trùng với ground-truth root cause hay không) giữa ba phương án: chỉ
dùng Rule Engine, chỉ dùng GraphRAG, và kết hợp cả hai qua module tổng hợp. Đánh giá được thực hiện
hai lần độc lập: một lần với \code{MockLLM} (xác định, không tốn chi phí, dùng để đánh giá số lượng
lớn 78 sự cố DejaVu), và một lần với LLM thật (Groq) trên tập viễn thông để kiểm tra hành vi trong
điều kiện thực tế.

\section{Kết quả trên tập viễn thông mô phỏng}

\begin{table}[H]
\centering
\caption{Kết quả trên 5 kịch bản viễn thông (MockLLM)}
\label{tab:eval-telecom}
\begin{tabular}{@{}lllll@{}}
\toprule
Kịch bản & Ground truth & Rule-only & GraphRAG-only & Combined \\
\midrule
SCN-01 & TRANS-CORE-1 & \checkmark & \checkmark & \checkmark \\
SCN-02 & RTR-EDGE-2   & \checkmark & \checkmark & \checkmark \\
SCN-03 & RTR-EDGE-1   & \checkmark & \checkmark & \checkmark \\
SCN-04 & GNB-A2       & \checkmark & \checkmark & \checkmark \\
SCN-05 & OLT-B1       & \checkmark & \checkmark & \checkmark \\
\midrule
\textbf{Độ chính xác} & & \textbf{100\% (5/5)} & \textbf{100\% (5/5)} & \textbf{100\% (5/5)} \\
\bottomrule
\end{tabular}
\end{table}

\section{Kết quả trên 78 sự cố thật DejaVu}

\begin{table}[H]
\centering
\caption{Kết quả trên 78 sự cố thật (MockLLM)}
\label{tab:eval-dejavu}
\begin{tabular}{@{}lc@{}}
\toprule
Phương án & Độ chính xác \\
\midrule
Rule-only     & 100\% (78/78) \\
GraphRAG-only & 100\% (78/78) \\
Combined      & 100\% (78/78) \\
\bottomrule
\end{tabular}
\end{table}

Kết quả 100\% trên cả hai tập với MockLLM cần được diễn giải cẩn trọng: vì toàn bộ sự cố DejaVu là
lỗi đơn-điểm, và luật~R2 (isolated device fault) được thiết kế chính xác cho mẫu hình này, kết quả
cao là hệ quả tự nhiên của việc thiết kế luật đúng bản chất bài toán, không phải bằng chứng cho
thấy hệ thống "học" được từ dữ liệu DejaVu (hệ thống không có bước huấn luyện). Giá trị của kết quả
này nằm ở chỗ: 4 luật và cơ chế MockLLM được xây dựng \emph{trước và độc lập} với việc quan sát dữ
liệu DejaVu, và vẫn generalize đúng sang một topology DAG đa-cha thật, phức tạp hơn hẳn tập viễn
thông tự tạo.

\section{Thử nghiệm với LLM thật và phát hiện về hiệu chỉnh độ tin cậy}
\label{sec:llm-thuc-nghiem}

Khi thay \code{MockLLM} bằng \code{GroqLLM} (mô hình \code{llama-3.3-70b-versatile}, gọi API thật)
trên 5 kịch bản viễn thông, GraphRAG-only đạt 4/5 (80\%): kịch bản SCN-05 (suy giảm KPI dần, không
có alarm cứng) bị LLM chẩn đoán sai thành \code{RTR-EDGE-2} thay vì \code{OLT-B1} đúng. Đáng chú ý
hơn, ở lần thử nghiệm đầu tiên (trước khi thêm cơ chế chiết khấu ở Mục~\ref{sec:combiner}), module
tổng hợp \textbf{cũng cho kết quả sai theo GraphRAG}, vì độ tin cậy LLM tự báo cáo cho câu trả lời
sai lại cao hơn độ tin cậy của ứng viên đúng từ Rule Engine (0.65).

Đây là một phát hiện thực nghiệm quan trọng, không phải lỗi cài đặt: LLM thật, khác với MockLLM
(vốn được lập trình để trả về độ tin cậy có hiệu chỉnh), không đảm bảo độ tin cậy tự báo cáo phản
ánh đúng xác suất chính xác thực tế -- một vấn đề đã được ghi nhận rộng rãi trong tài liệu về hiệu
chỉnh độ tin cậy (confidence calibration) của LLM~\cite{guo2017calibration}. Sau khi bổ sung cơ chế
chiết khấu ứng viên GraphRAG không được luật xác nhận (Đoạn mã~\ref{lst:discount}), kết quả
\emph{combined} khôi phục về 5/5 (100\%) mà không làm giảm độ chính xác ở bất kỳ trường hợp nào
khác (đã chạy lại toàn bộ 5 kịch bản viễn thông và 78 sự cố DejaVu để xác nhận). Kết quả này minh
hoạ cụ thể, bằng số liệu thực nghiệm, cho luận điểm lý thuyết đã nêu ở Mục~\ref{sec:crossval-theory}
và trong đề cương ban đầu: suy diễn dựa luật đóng vai trò kiểm chứng chéo không thể thay thế cho
suy luận của LLM, ngay cả khi LLM đã có ngữ cảnh được truy xuất chính xác từ Knowledge Graph.

\section{Phân tích lỗi có hệ thống trên mẫu mở rộng (Error Analysis)}
\label{sec:error-analysis}

Kết quả ở Mục~\ref{sec:llm-thuc-nghiem} đặt ra một câu hỏi: liệu sai lầm của GroqLLM ở SCN-05 là
ngẫu nhiên, hay phản ánh một xu hướng suy luận sai có hệ thống? Vì độ chính xác top-1 đã bão hoà ở
100\% trên cả hai tập dữ liệu khi dùng MockLLM, chạy toàn bộ 78 sự cố DejaVu qua LLM thật sẽ tốn
thời gian/quota API mà không chắc sinh thêm thông tin định lượng mới -- một đánh đổi (trade-off)
kinh điển khi một độ đo đã chạm trần (ceiling effect): thay vì thêm số liệu tổng hợp, giá trị nằm ở
phân tích định tính các trường hợp sai. Một mẫu 20 sự cố được chọn có chủ đích để tối đa hoá khả
năng LLM mắc lỗi: toàn bộ 5 kịch bản viễn thông, cộng 15 sự cố DejaVu được chọn theo \textbf{độ khó
cấu trúc} -- số thiết bị cha (\code{DEPENDS\_ON}) của thiết bị bị cảnh báo, vì thiết bị càng nhiều
cha thì càng có nhiều "nghi phạm giả" khiến suy luận dựa trên ngữ cảnh dễ đi lạc hướng. Cụ thể: 12
sự cố trên bốn thiết bị "khó nhất" (\code{docker\_001} đến \code{docker\_004}, mỗi thiết bị có 5
thiết bị cha), phủ đủ tổ hợp (thiết bị, loại lỗi) có thể xảy ra trên nhóm này; và 3 sự cố trên nhóm
thiết bị có 2 thiết bị cha, để có điểm so sánh phổ độ khó.

Bảng~\ref{tab:error-analysis} so sánh MockLLM và GroqLLM trên đúng mẫu 20 sự cố này.

\begin{table}[H]
\centering
\caption{So sánh MockLLM và GroqLLM (LLM thật) trên mẫu 20 sự cố chọn lọc theo độ khó}
\label{tab:error-analysis}
\begin{tabular}{@{}lcc@{}}
\toprule
Phương án & MockLLM & GroqLLM (LLM thật) \\
\midrule
Rule-only     & 100\% (20/20) & 100\% (20/20) \\
GraphRAG-only & 100\% (20/20) & \textbf{45\% (9/20)} \\
Combined      & 100\% (20/20) & \textbf{100\% (20/20)} \\
\bottomrule
\end{tabular}
\end{table}

Ba quan sát quan trọng:

\begin{enumerate}[label=(\arabic*)]
    \item Rule Engine không đổi (100\% cả hai cột) -- đúng như kỳ vọng, vì nó không phụ thuộc LLM.
    \item GraphRAG-only \textbf{sụt từ 100\% xuống 45\%} khi chuyển từ MockLLM sang LLM thật, trên
    đúng mẫu được chọn để "làm khó" nó -- một khoảng cách lớn hơn nhiều so với con số 80\% quan sát
    được ở SCN-05 đơn lẻ (Mục~\ref{sec:llm-thuc-nghiem}), cho thấy đây không phải hiện tượng hiếm.
    \item Combined vẫn giữ nguyên \textbf{100\%} -- toàn bộ 11/11 trường hợp GraphRAG trả lời sai
    đều được module tổng hợp sửa đúng nhờ cơ chế chiết khấu (Mục~\ref{sec:combiner}).
\end{enumerate}

\paragraph{Một khuôn mẫu lỗi duy nhất, không phải nhiễu ngẫu nhiên.} Xem xét chi tiết cả 11 trường
hợp GraphRAG trả lời sai (Bảng~\ref{tab:error-cases}) cho thấy chúng đều theo \textbf{đúng một khuôn
mẫu}: LLM luôn chọn một \emph{thiết bị cha} (\code{db\_XXX} hoặc \code{os\_XXX}) của thiết bị thực
sự bị lỗi làm nguyên nhân gốc rễ, với lý do "thiết bị cha này được nhiều thiết bị khác cùng phụ
thuộc, nên có thể là nguyên nhân chung" -- trong khi thiết bị cha đó \textbf{không hề có alarm hay
bằng chứng bất thường nào} trong ngữ cảnh được cung cấp. Đây là một dạng suy luận quy nạp thiếu căn
cứ (unwarranted generalization): LLM đánh đồng "được nhiều thiết bị phụ thuộc" với "có khả năng là
nguyên nhân" -- một suy luận hợp lý \emph{trong trường hợp cascade thật} (đúng như luật R3 của Rule
Engine mã hoá tường minh, Bảng~\ref{tab:rules}) nhưng sai trong các trường hợp này, vì toàn bộ 78 sự
cố DejaVu đều là lỗi đơn-thiết-bị (Mục~\ref{ch:data}) -- không có thiết bị nào khác thực sự bị ảnh
hưởng để biện minh cho suy luận đó.

\begin{table}[H]
\centering
\small
\caption{Toàn bộ 11 trường hợp GraphRAG (LLM thật) trả lời sai trong mẫu 20 sự cố}
\label{tab:error-cases}
\begin{tabular}{@{}p{1.7cm}p{1.7cm}p{1.7cm}p{7.0cm}@{}}
\toprule
ID & Ground truth & GraphRAG đoán & Tóm tắt lý do LLM đưa ra (đều theo cùng một khuôn mẫu) \\
\midrule
SCN-05 & OLT-B1 & RTR-EDGE-2 & Không có alarm, chỉ KPI bất thường ở OLT-B1; nhầm sang case lịch sử liên quan router \\
DEJAVU-002 & docker\_002 & db\_009 & "db\_009 là thiết bị nhiều docker khác cùng phụ thuộc" \\
DEJAVU-003 & docker\_001 & docker\_007 & Nhầm sang một thiết bị ngang hàng không liên quan trực tiếp \\
DEJAVU-004 & docker\_004 & db\_007 & "db\_007 là thiết bị chung của nhiều docker" \\
DEJAVU-006 & docker\_002 & db\_009 & "db\_009 là thiết bị chung cho nhiều docker" \\
DEJAVU-008 & docker\_003 & os\_019 & "os\_019 là thiết bị có thể ảnh hưởng đến docker\_003" \\
DEJAVU-013 & docker\_001 & db\_003 & Suy luận bắc cầu qua một thiết bị cha không liên quan trực tiếp \\
DEJAVU-029 & docker\_002 & db\_009 & "db\_009 được nhiều docker khác phụ thuộc" \\
DEJAVU-032 & docker\_004 & db\_007 & "db\_007 là thiết bị chung, không có thông tin nên nghi ngờ nó" \\
DEJAVU-010 & docker\_007 & db\_003 & "db\_003 cũng là thiết bị phụ thuộc của các docker khác" \\
DEJAVU-014 & docker\_005 & db\_003 & "các thiết bị khác phụ thuộc db\_003 cũng có thể bị ảnh hưởng" \\
\bottomrule
\end{tabular}
\end{table}

Phát hiện này có ý nghĩa thực tiễn rõ ràng cho một hệ thống hỗ trợ NOC: nếu triển khai GraphRAG đơn
lẻ (không có Rule Engine), hệ thống sẽ có xu hướng \textbf{hệ thống} (không ngẫu nhiên) đổ oan cho
hạ tầng dùng chung (database, host) mỗi khi một dịch vụ đơn lẻ gặp sự cố -- một sai lầm chẩn đoán
tốn kém trong vận hành thực tế (điều tra nhầm hướng, có thể ảnh hưởng dây chuyền đến các dịch vụ
khác đang chạy tốt trên cùng hạ tầng dùng chung). Việc Rule Engine sửa đúng toàn bộ 11/11 trường hợp
này không phải may mắn: luật R2 (isolated device fault, Bảng~\ref{tab:rules}) kiểm tra chính xác
điều kiện LLM đã bỏ qua -- liệu thiết bị cha có bằng chứng bất thường hay không, chứ không suy luận
từ cấu trúc phụ thuộc đơn thuần. Toàn bộ dữ liệu thô (bao gồm nguyên văn giải thích của LLM cho từng
trường hợp) được lưu tại \code{report/real\_llm\_sample\_results.json}, sinh ra bởi
\code{src/evaluation/evaluate\_real\_llm\_sample.py}.

\section{Kiểm thử phần mềm}

Bộ kiểm thử tự động (\code{pytest}) gồm 14 test case, bao trùm: logic từng luật của Rule Engine,
công thức kết hợp certainty factor, bộ phân tích ngữ cảnh (\code{context\_builder}), suy luận của
MockLLM (case-based matching và suy luận cấu trúc), cơ chế chiết khấu ở module tổng hợp, và tính
toàn vẹn của dữ liệu DejaVu đã tải về. Toàn bộ test không phụ thuộc Neo4j đang chạy hay không (dữ
kiện được dựng trực tiếp bằng cấu trúc Python), giúp chạy nhanh (dưới 0.1 giây) và có thể tích hợp
vào quy trình kiểm tra liên tục.

% ==============================================================
\chapter{Kết luận và hướng phát triển}

\section{Kết luận}

Báo cáo đã trình bày một hệ thống hoàn chỉnh, cài đặt được và kiểm chứng được, minh hoạ ba kỹ thuật
biểu diễn và suy diễn tri thức của học phần trên cùng một bài toán thực tế: Knowledge Graph để biểu
diễn tường minh cấu trúc phụ thuộc, suy diễn dựa luật theo mô hình production system với cơ chế
certainty factor kiểu MYCIN, và GraphRAG để khai thác LLM một cách có kiểm soát. Hai đóng góp thực
nghiệm đáng chú ý ngoài phạm vi tối thiểu của đề cương ban đầu là: (1) tổng quát hoá kiến trúc sang
đồ thị phụ thuộc đa-cha và kiểm chứng độc lập trên 78 sự cố hạ tầng thật (NetManAIOps/DejaVu), và
(2) phát hiện và khắc phục vấn đề hiệu chỉnh độ tin cậy khi tích hợp LLM thật, cho thấy giá trị cụ
thể, đo lường được của cơ chế kiểm chứng chéo giữa luật và LLM.

\section{Hướng phát triển}

\begin{itemize}
    \item Mở rộng bộ luật để xử lý các mẫu hình lan truyền phức tạp hơn (nhiều tầng cascade đồng
    thời, sự cố gián đoạn/không liên tục).
    \item Đánh giá thêm với các LLM khác (mô hình lớn hơn, hoặc mô hình được fine-tune riêng cho
    domain vận hành mạng) để so sánh mức độ hiệu chỉnh độ tin cậy.
    \item Xây dựng cơ chế học/cập nhật trọng số luật từ phản hồi của kỹ sư vận hành (human-in-the
    -loop), thay vì trọng số cố định như hiện tại.
    \item Nếu có điều kiện tiếp cận dữ liệu vận hành thật của một nhà mạng (qua thoả thuận bảo mật),
    kiểm chứng lại toàn bộ phương pháp trên đúng miền viễn thông với dữ liệu thật.
\end{itemize}

% ==============================================================
\begin{thebibliography}{99}

\bibitem{shortliffe1975mycin}
E.~H. Shortliffe and B.~G. Buchanan,
\emph{A model of inexact reasoning in medicine},
Mathematical Biosciences, 23(3-4):351--379, 1975.

\bibitem{guo2017calibration}
C.~Guo, G.~Pleiss, Y.~Sun, K.~Q. Weinberger,
\emph{On Calibration of Modern Neural Networks},
Proceedings of the 34th International Conference on Machine Learning (ICML), 2017.

\bibitem{netrca2022}
Signal Processing Society, IEEE,
\emph{Root Cause Analysis for Wireless Network Faults Localization (NetRCA) -- ICASSP 2022 AIOps
Challenge}.
\url{https://signalprocessingsociety.org/publications-resources/data-challenges/root-cause-analysis-wireless-network-faults-localization}

\bibitem{telelogs2025}
NetOp Team, Huawei Paris Research Center,
\emph{TeleLogs: A Benchmark for Large Language Models in Telecom Data Analysis},
Hugging Face Datasets, \url{https://huggingface.co/datasets/netop/TeleLogs}, 2025.

\bibitem{dejavu2022}
M.~Li, Z.~Li, K.~Yin, X.~Nie, W.~Zhang, K.~Sui, D.~Pei,
\emph{Actionable and Interpretable Fault Localization for Recurring Failures in Online Service
Systems},
Proceedings of the 30th ACM Joint European Software Engineering Conference and Symposium on the
Foundations of Software Engineering (ESEC/FSE), 2022.
Dữ liệu: Zenodo, DOI: 10.5281/zenodo.6955909, giấy phép CC-BY~4.0.

\bibitem{neo4j}
Neo4j, Inc., \emph{Neo4j Graph Database Documentation}, \url{https://neo4j.com/docs/}.

\end{thebibliography}

\end{document}
